\section{Runtime Verification Và Service Mesh}

\subsection{Platform Readiness}

Platform stack gồm PostgreSQL, Redis, Kafka, Elasticsearch và Keycloak. Các service
YAS ở \texttt{dev} và \texttt{staging} dùng chung platform nhưng tách database theo
environment. Evidence hiện có xác nhận node, storage class, PVC và database đã sẵn sàng.

\begin{figure}[H]
\centering
\reportimg{img/platform/02_gcp_vm_ip_addr.png}
\caption{External IP của GCP VM dùng cho NodePort demo}
\end{figure}

\subsection{External Access}

Runtime access chính dùng Nginx ingress NodePort \texttt{30846} thông qua cơ chế phân phối Multi-Host Ingress, định tuyến bằng trường Host header. Các tên miền dùng chung cũ (\texttt{yas.<env>.local}) đã được loại bỏ hoàn toàn để phân tách luồng nghiệp vụ giữa storefront, backoffice, swagger và identity. Developer hoặc giảng viên thêm host mapping vào tệp hosts:

\begin{lstlisting}[language=bash, caption={Cấu hình tệp hosts phân giải tên miền mới}]
34.124.212.254 storefront.dev.yas.local.com backoffice.dev.yas.local.com swagger.dev.yas.local.com identity.dev.yas.local.com
34.124.212.254 storefront.staging.yas.local.com backoffice.staging.yas.local.com swagger.staging.yas.local.com identity.staging.yas.local.com
34.124.212.254 storefront.developer.yas.local.com backoffice.developer.yas.local.com swagger.developer.yas.local.com identity.developer.yas.local.com
\end{lstlisting}

Storefront dev/staging đã được smoke test qua \texttt{curl} theo tên miền \path{storefront.<env>.yas.local.com}: HTML trả 200, product
API trả seeded iPhone data, media thumbnail trả \texttt{image/jpeg}, và Keycloak
login/register render đúng host cùng NodePort.

\begin{figure}[H]
\centering
\reportimg{img/runtime/17_app_reachable_dev.png}
\caption{Storefront dev reachable qua NodePort}
\end{figure}

\subsection{Service Mesh Scope}

Mesh được áp dụng cho \texttt{dev} và \texttt{staging}, không chỉ một namespace demo
tách biệt. Các pod ứng dụng có workload container và \texttt{istio-proxy}, do đó trạng
thái READY kỳ vọng là \texttt{2/2}. Mesh policy gồm:

\begin{itemize}
    \item PeerAuthentication cho STRICT mTLS ở backend path.
    \item DestinationRule cho traffic mesh.
    \item VirtualService cho retry/fault evidence.
    \item AuthorizationPolicy theo service identity để có allow/deny rõ ràng.
\end{itemize}

\begin{table}[H]
\centering
\caption{Service mesh requirement và implementation}
\begin{tabular}{|p{3.2cm}|p{5.1cm}|p{5.2cm}|}
\hline
\textbf{Yêu cầu mesh} & \textbf{Triển khai trong CD repo} & \textbf{Evidence cần trình bày} \\
\hline
mTLS & \path{overlays/dev/istio/mtls.yaml} và \path{overlays/staging/istio/mtls.yaml} dùng \texttt{PeerAuthentication} STRICT và DestinationRule \texttt{ISTIO\_MUTUAL}. & Screenshot pod \texttt{READY 2/2}, manifest mTLS, Kiali security lock nếu có. \\
\hline
Retry & \path{virtual-services.yaml} khai báo retry cho critical services với \texttt{retryOn: 5xx,connect-failure,refused-stream}. & Manifest VirtualService và log/curl test retry. \\
\hline
Authorization & \path{authorization-policies.yaml} giới hạn service-to-service bằng service account/SPIFFE principal. & Curl allow/deny: \texttt{tax -> location} allowed, \texttt{cart -> location} denied 403. \\
\hline
Topology & Kiali quan sát graph traffic dev/staging. & Screenshot Kiali topology hiển thị storefront/BFF/backend path. \\
\hline
\end{tabular}
\end{table}

\begin{figure}[H]
\centering
\reportimg{img/mesh/mtls_strict_pods_2of2.png}
\caption{Trạng thái các Pod trong namespace dev hoạt động với 2/2 Container (chứa Envoy sidecar)}
\end{figure}

\subsection{AuthorizationPolicy Evidence}

Kịch bản allow/deny dùng đường gọi \texttt{tax -> location}:

\begin{itemize}
    \item Gọi từ \texttt{tax} sang \texttt{location}: request đi qua mesh và đến application layer.
    \item Gọi từ \texttt{cart} sang \texttt{location}: bị chặn HTTP 403 bởi Istio AuthorizationPolicy.
\end{itemize}

HTTP 403 trong case denied là bằng chứng policy hoạt động ở mesh layer. HTTP 500 ở
case allowed vẫn có giá trị vì nó chứng minh request đã qua được mesh và chạm vào
Spring Boot application, lỗi còn lại thuộc application endpoint.

\begin{figure}[H]
\centering
\reportimg{img/mesh/authorization_allow_tax_location.png}
\caption{Bằng chứng kết nối thành công (Allowed) từ tax service sang location service}
\end{figure}

\begin{figure}[H]
\centering
\reportimg{img/mesh/authorization_deny_cart_location.png}
\caption{Bằng chứng kết nối bị chặn (Denied HTTP 403) từ cart service sang location service bởi Envoy RBAC}
\end{figure}

\begin{figure}[H]
\centering
\reportimg{img/mesh/retry_policy_evidence.png}
\caption{Minh chứng log Envoy thực hiện cơ chế Retry tự động (3 lần) khi service đích gặp lỗi}
\end{figure}

\subsection{Kiali Topology}

Kiali được dùng để quan sát topology service mesh. Nếu chưa có ảnh cuối cùng trong
\path{latex-report/img/mesh/}, báo cáo cần để placeholder rõ: chụp Kiali topology
cho \texttt{dev} hoặc \texttt{staging}, hiển thị luồng storefront/BFF/backend và
sidecar traffic.

\begin{figure}[H]
\centering
\reportimg{img/mesh/kiali_topology_dev_staging.png}
\caption{Sơ đồ Graph Topology các service trong mesh quan sát được trên giao diện Kiali}
\end{figure}
